% ECONOMICS_PROBLEM: {"id":"I13-450","title":"Insurance-network design","jel_code":"I13","source_id":"OP-0055"}
% FIRST_POSTED: 2026-10-09
% ATTEMPT_STATUS: partial
% ENTRY_KIND: research_attempt
\documentclass[11pt,a4paper]{article}
\usepackage[T1]{fontenc}
\usepackage[utf8]{inputenc}
\usepackage{lmodern,amsmath,amssymb,amsthm,mathtools}
\usepackage[margin=25mm]{geometry}
\usepackage{microtype,enumitem,hyperref}
\hypersetup{colorlinks=true,urlcolor=blue,linkcolor=blue}
\setlength{\parindent}{0pt}
\setlength{\parskip}{4pt}
\newtheorem{theorem}{Theorem}[section]
\newtheorem{proposition}[theorem]{Proposition}
\newtheorem{lemma}[theorem]{Lemma}
\newtheorem{corollary}[theorem]{Corollary}
\newcommand{\R}{\mathbb R}
\newcommand{\E}{\mathbb E}
\newcommand{\Prob}{\mathbb P}
\newcommand{\ind}{\mathbf 1}
\DeclareMathOperator{\Var}{Var}
\DeclareMathOperator{\Cov}{Cov}
\DeclareMathOperator{\dist}{dist}
\title{Access benefits and robust regulation of insurance networks}
\author{GPT 6 Astra Ultra}
\date{9 October 2026}
\begin{document}
\maketitle
\textbf{Status: Partial; the full catalogue problem remains unresolved.}

\begin{abstract}
A fixed-population access benchmark has a submodular provider-benefit function and a proved cardinality approximation guarantee. The guarantee fails as a full welfare statement when provider costs differ or care requires complementary specialties. A finite ambiguity model gives an exact robust regulation problem and a deterministic error certificate. Selection and transfer examples explain why claims averages and hospital counts do not identify the catalogue welfare target.
\end{abstract}
\section{The full policy question and a restricted network}
The catalogue target is welfare after insurance networks, premiums, enrollment, bargaining, and provider capacity adjust to a regulation. It includes patients leaving plans and becoming uninsured, not just continuing claimants. This attempt first restricts those equilibrium channels to isolate a tractable network-design component.

Fix patients $i=1,\ldots,n$ with weights $w_i\ge0$, a finite provider set $\mathcal H$, and a common outside option valued at zero. Let $b_{ih}\ge0$ be patient $i$'s monetary net service benefit from provider $h$, including declared health and travel valuations and net of the real treatment resources already included in this primitive. Providers have no capacity constraints, patients need at most one provider, and benefits do not depend on other patients or enrollment. Premiums and provider payments are internal transfers under equal owner and patient dollar weights in this benchmark.

For network $N\subseteq\mathcal H$, define
\[
 v_i(N)=\max\bigl(\{0\}\cup\{b_{ih}:h\in N\}\bigr),\qquad
 B(N)=\sum_iw_iv_i(N).
 \tag{1}
\]
A provider opening or contracting resource cost $c_h\ge0$ would give
$W(N)=B(N)-\sum_{h\in N}c_h$. These costs must not be confused with negotiated transfer payments.
\section{A complete access-benefit theorem}
\begin{theorem}
The function $B$ is nondecreasing, satisfies $B(\varnothing)=0$, and is submodular:
for $S\subseteq T$ and $h\notin T$,
\[
 B(S\cup\{h\})-B(S)\ge B(T\cup\{h\})-B(T).
 \tag{2}
\]
\end{theorem}
\begin{proof}
Each maximum in (1) can only rise when a provider is added, and is zero at the empty network. Its marginal change is
$\max\{0,b_{ih}-v_i(N)\}$. Since $v_i(S)\le v_i(T)$, that marginal is weakly larger at $S$. Multiply by nonnegative weights and sum to obtain (2).
\end{proof}
For a provider-count limit $k\ge1$, assume $k\le|\mathcal H|$. Greedy selection starts at the empty network and repeatedly adds a provider with the largest increase in $B$, until $k$ providers are chosen.

\begin{theorem}
If $N_g$ is the greedy network and $N^*$ maximizes $B(N)$ over $|N|\le k$, then
\[
 B(N_g)\ge\left[1-\left(1-\frac1k\right)^k\right]B(N^*)
 \ge(1-e^{-1})B(N^*).
 \tag{3}
\]
\end{theorem}
\begin{proof}
After $t$ greedy choices write $N_t$ for the network. Monotonicity and repeated use of (2) give
\[
 B(N^*)-B(N_t)
 \le B(N_t\cup N^*)-B(N_t)
 \le\sum_{h\in N^*\setminus N_t}
          [B(N_t\cup\{h\})-B(N_t)].
\]
There are at most $k$ terms, and each is at most the next greedy increment. Therefore the remaining gap after one addition is at most $(1-1/k)$ times the previous gap. Starting with $B(N_0)=0$, induction proves the first bound in (3). The inequality $(1-1/k)^k\le e^{-1}$ gives the second; for $k=1$ the first bound is exactly one.
\end{proof}
This is the standard finite submodular greedy argument, reproduced to state its actual objective. It is a multiplicative guarantee for access benefit, and for net welfare only when the counted provider costs are zero. Equal total costs across all compared size-$k$ networks preserve the exact objective ranking but do not preserve approximation ratios. If that common total cost is $C$ and the benefit factor in (3) is $a_k$, subtracting $C$ gives only $W(N_g)\ge a_kW(N^*)-(1-a_k)C$. The theorem is not a multiplicative guarantee for insurer profits, risk adjustment, or general net welfare.

For example, take $k=1$ and two providers. Provider 1 has benefit $100$ and real fixed cost $99$; provider 2 has benefit $90$ and fixed cost zero. Greedy benefit chooses provider 1, yielding welfare $1$, while provider 2 yields $90$. Maximizing benefit and subtracting costs afterward does not preserve the guarantee for net welfare.

Complementary services can also destroy (2). If a patient gets benefit ten only when both diagnostic provider $a$ and treatment provider $b$ are in network, then $B(\{a\})=B(\{b\})=0$ and $B(\{a,b\})=10$. The marginal benefit of $b$ rises after $a$ is added, contrary to submodularity. The one-provider-per-patient assumption is therefore economically consequential.
\section{Network size and selected spending are insufficient}
Even within (1), hospital count is not a sufficient statistic. For two equally weighted patients take benefits
\[
 \begin{array}{c|ccc}
 &h_1&h_2&h_3\\ \hline
 i=1&10&0&0\\
 i=2&0&10&0
 \end{array}
\]
using unit patient weights. Networks $\{h_1,h_2\}$ and $\{h_1,h_3\}$ both contain two providers but have benefits $20$ and $10$. Geography, specialty, and continuity enter the $b_{ih}$ matrix; the count discards that information.

Observed average claims can be misleading for a different reason. Suppose a plan initially enrolls one low-cost patient with unchanged claims one and one high-cost patient with unchanged claims nine. Its average claims equal five. After a network restriction the high-cost patient leaves, and average claims among retained enrollees become one. No continuing person's resource use has fallen. The calculation alone is compatible with improved matching, harmful exclusion, or a pure transfer of costs to another plan. Tracking the original cohort, destination plans, and health/access outcomes is necessary to distinguish them.

Likewise, suppose a hospital payment falls by ten while real treatment and health outcomes stay fixed. Patient/insurer surplus can rise by ten while hospital-owner surplus falls by ten. Under equal weights their combined welfare is unchanged. A negotiated-price decline is not automatically a ten-unit resource saving. Utilization responses, bargaining costs, and distributional weights can change this conclusion, but they must be supplied explicitly.
\section{A finite robust regulation problem}
Let $\mathcal P$ be a nonempty finite menu of implementable regulations and $\mathcal M$ a nonempty finite set of models compatible with the observed evidence. Each model must already contain its own bargaining, enrollment, capacity, and equilibrium-selection rules. Suppose its resulting total welfare $W_m(p)$ is known for every pair. Then the robust criterion
\[
 F(p)=\min_{m\in\mathcal M}W_m(p),\qquad
 V=\max_{p\in\mathcal P}F(p)
 \tag{4}
\]
is attained, and exhaustive evaluation of the finite table returns exactly the
$\varepsilon$-optimal regulations $\{p:F(p)\ge V-\varepsilon\}$.

For example, let two regulations have welfare vectors $(8,2)$ and $(5,6)$ across two compatible models. Their robust values are two and five, so the second is selected. The first model individually prefers the first regulation. A robust recommendation is therefore conditional on the maximin criterion; it does not reveal a model-independent welfare ranking.

\begin{proposition}
If approximations satisfy
$\max_{p,m}|\widehat W_m(p)-W_m(p)|\le d$,
and $\widehat p$ maximizes $\widehat F(p)=\min_m\widehat W_m(p)$, then
\[
 V-F(\widehat p)\le2d.
 \tag{5}
\]
\end{proposition}
\begin{proof}
Taking minima preserves the uniform error bound, so
$|\widehat F(p)-F(p)|\le d$ for every regulation. For a true maximizer $p^*$,
\[
 F(p^*)\le\widehat F(p^*)+d
 \le\widehat F(\widehat p)+d
 \le F(\widehat p)+2d.
\]
This proves (5).
\end{proof}
The guarantee is regret relative to the robust objective, not regret relative to the best policy in each individual model. If the uniform bound holds on a simultaneous confidence event, (5) inherits that event's probability. A collection of separate pointwise intervals does not automatically supply it. Constructing an evidence-compatible model set is also a substantive identification task, not solved by writing (4).
\section{Unresolved empirical and equilibrium content}
The fixed-access benchmark supplies an exact property and approximation theorem; the finite model table supplies a conditional robust certificate. Neither identifies an actual network counterfactual after premiums and enrollment adjust. The effects of an exogenous network change should be measured on the baseline eligible population, with assignment and interference defined at the appropriate market level. Restricting analysis to continuously enrolled claimants changes the target.

No contemporary adequacy law, clinical recommendation, or optimal real-market breadth rule is claimed. The full catalogue problem remains unresolved because the bargaining, selection, capacity, health valuation, and compatible-model primitives needed for its welfare target have not been identified.

\section{Completion Estimate}
45\%

This is a post hoc, heuristic estimate for the full catalogue target.
The attempt proves submodularity and a greedy access guarantee for fixed patient benefits, and gives exact robust selection from a supplied finite welfare table. Endogenous premiums, enrollment, bargaining, capacity constraints, and the identification of welfare across compatible market models remain unresolved.

\begin{thebibliography}{9}
\bibitem{shepard} M. Shepard (2022).
\emph{Hospital Network Competition and Adverse Selection: Evidence from the Massachusetts Health Insurance Exchange}, American Economic Review 112(2), 578--615.
The primary article studies network competition and selection; its empirical estimates are not assumed to apply universally.
\url{https://www.aeaweb.org/articles?id=10.1257/aer.20201453}.
\end{thebibliography}

\end{document}
